\section{Control de varianza y mejora de la eficiencia de muestreo}
\subsection{GAE y temporal smoothing}
Generalized Advantage Estimation (GAE) hace una transición suave entre reward-to-go y TD error:
$$\hat{A}_t^{GAE(\gamma,\lambda)} = \sum_{l=0}^{\infty}(\gamma\lambda)^l \delta_{t+l},$$
donde $\delta_t = r_t + \gamma V(s_{t+1}) - V(s_t)$. Ajustar $\lambda$ permite hacer un trade-off entre bias y variance.

\begin{importantbox}{La lógica de ajuste de GAE}
\begin{itemize}
    \item $\lambda=1$ corresponde al Monte Carlo completo (alta variance / baja bias);
    \item $\lambda=0$ corresponde a TD(0) (baja variance / alta bias);
    \item En la práctica, tomar $\lambda \in [0.9,0.98]$ suele lograr el mejor compromiso.
\end{itemize}
\end{importantbox}

\subsection{Regularización de Entropy y Trust Region}
Para fomentar la exploración, se agrega entropy regularization al loss:
$$\tilde{J}_{ent}(\theta) = \tilde{J}(\theta) + \beta \mathbb{E}[-\log \pi_\theta(a\mid s)].$$
El Entropy term evita que la política se determine prematuramente y caiga en un suboptimal mode. Otra práctica común es limitar la KL divergence, formando un trust region update (como en TRPO):
$$\max_\theta \mathbb{E}[\log \pi_\theta(a\mid s) \hat{A}] \quad \text{s.t. } D_{KL}(\pi_{\theta_{old}} \,\|\, \pi_\theta) \le \delta.$$

\begin{knowledgebox}{El equilibrio entre exploración y estabilidad}
\begin{itemize}
    \item Entropy regularization ofrece un soft constraint que fomenta una distribución de probabilidad más plana;
    \item El KL constraint limita explícitamente el tamaño de paso mediante el trust region, evitando el policy collapse;
    \item Una combinación típica es PPO (clipped objective) + entropy bonus.
\end{itemize}
\end{knowledgebox}

\subsection{Resumen de la sección}
GAE, entropy boost y trust region son técnicas habituales para mejorar la estabilidad y la eficiencia de muestreo en el entrenamiento real; mantienen la variance dentro de un rango manejable y conservan sufficient exploration.

